\chapter{Capitolo 1 - Crisi dell'ontologia estensionale e problema della validità}

\section{1. Problema}

La teoria delle categorie classica ha prodotto uno dei linguaggi più potenti della matematica contemporanea. Il suo punto di forza è noto: non descrive gli oggetti ``in sé'', ma le loro relazioni strutturali e le regole di composizione che rendono trasferibili risultati da un dominio all'altro. Questa astrazione ha reso possibile una unificazione profonda tra algebra, topologia, logica, geometria e informatica teorica.

Tuttavia, lo stesso gesto che rende la teoria straordinariamente generale introduce una domanda che qui prendiamo come problema fondativo: \textbf{quando una struttura sia formalmente ben composta sia ontologicamente valida come struttura reale?}

In sintesi, il problema non riguarda la correttezza sintattica della category theory, che resta intatta, ma la sua neutralità rispetto a tre aspetti che diventano decisivi quando si passa da domini puramente formali a sistemi vivi, storici, cognitivi o socio-tecnici:

\begin{enumerate}
\def\labelenumi{\arabic{enumi}.}
\tightlist
\item
  la differenza tra composizione possibile e composizione feconda;
\item
  la differenza tra equivalenza formale e equivalenza funzionale;
\item
  la differenza tra universalità strutturale e capacità generativa.
\end{enumerate}

Se due costruzioni sono isomorfe nel senso classico, la teoria ha pieno titolo per trattarle come ``le stesse'' ai fini della struttura. Ma in molti domini empirici questa identificazione è troppo forte: sistemi isomorfi in una proiezione possono divergere radicalmente nella stabilità, nella robustezza, nella capacità emergente e nella traiettoria dinamica.

Il punto non è negare la validità del classico. Il punto è riconoscere un \textbf{deficit di discriminazione ontologica} quando il linguaggio classico viene usato come grammatica di sistemi reali complessi.

Da qui la tesi metodologica del capitolo: occorre distinguere in modo esplicito tra:

\begin{itemize}
\tightlist
\item
  validità sintattica (ben-tipizzazione, commutatività, universalità, equivalenza strutturale);
\item
  validità ordinativa (coerenza interna non degenerativa + emergenza funzionale non nulla in contesto osservativo esplicito).
\end{itemize}

Questa distinzione apre il campo della Ordinative Category Theory (OCT): un'estensione conservativa della sintassi categoriale classica tramite uno strato di validazione ontologica.

\subsection{1.1 Posizionamento storico minimo del problema}

Il problema non nasce in astratto. Ogni grande formalismo passa attraverso una fase in cui la sua forza tecnica produce un effetto collaterale: tende a essere usato oltre il proprio mandato originario.

Per la teoria delle categorie, il mandato originario è chiarissimo: descrivere invarianti di struttura e trasformazioni naturali tra strutture. In questo mandato, il formalismo è difficilmente superabile. Ma quando lo stesso strumento viene trasposto in domini dove la differenza tra ``funziona in teoria'' e ``regge nel reale'' è decisiva, emerge un varco metodologico.

Questo varco è visibile in molti contesti contemporanei:

\begin{enumerate}
\def\labelenumi{\arabic{enumi}.}
\tightlist
\item
  in architetture computazionali formalmente corrette ma instabili su distribuzioni reali;
\item
  in modelli sociali coerenti in simulazione ma non robusti in campo;
\item
  in pipeline semantiche con ottima composizione locale ma collasso globale di senso.
\end{enumerate}

Nessuno di questi esempi invalida la categoria classica. Mostra solo che la domanda scientifica si è ampliata:

\begin{itemize}
\tightlist
\item
  non solo ``qual è la forma?'',
\item
  ma anche ``quale forma resta operativa sotto stress di realtà?''.
\end{itemize}

OCT assume questo ampliamento come punto di avvio e non come critica ideologica.

\subsection{1.2 Perché ``estensionale'' qui è una nozione tecnica}

Nel capitolo, ``estensionale'' non è usato come etichetta polemica. Indica una postura in cui l'identità viene principalmente stabilità da criteri di equivalenza/sostituibilità strutturale.

Questa postura ha tre virtu:

\begin{enumerate}
\def\labelenumi{\arabic{enumi}.}
\tightlist
\item
  pulizia formale;
\item
  trasferibilità dei risultati;
\item
  forte economia concettuale.
\end{enumerate}

Ma in sistemi ad alta dipendenza relazionale ha anche un costo:

\begin{enumerate}
\def\labelenumi{\arabic{enumi}.}
\tightlist
\item
  tende a comprimere differenze di stato interno che diventano decisive nel tempo;
\item
  rende opaco il passaggio tra ``possibile'' e ``sostenibile'';
\item
  non dispone, per progetto, di un criterio interno di non-degenerazione.
\end{enumerate}

L'estensione ordinativa non elimina il regime estensionale. Lo completa con un regime di valutazione che esplicita la dimensione di persistenza funzionale.

\subsection{1.3 Cosa significa ``validità'' in un testo OCT}

Nel resto del libro, la parola ``validità'' è sempre usata in senso stratificato:

\begin{enumerate}
\def\labelenumi{\arabic{enumi}.}
\tightlist
\item
  \textbf{validità formale}: la costruzione è corretta nel suo sistema di regole;
\item
  \textbf{validità ordinativa}: la costruzione preserva coerenza e non collassa la funzione emergente nel contesto dichiarato;
\item
  \textbf{validità scientifica}: il claim associato alla costruzione è sostenuto da evidenza classificata e protocollo riproducibile.
\end{enumerate}

La mancata distinzione tra questi tre livelli è la prima fonte di confusione nella discussione pubblica delle teorie complesse. Questo capitolo imposta la distinzione come regola permanente di lettura.

\subsection{1.4 Percorsi di lettura (ingresso rapido e percorso completo)}

Per ridurre la soglia di ingresso senza indebolire il rigore, il libro adotta due percorsi ufficiali.

Percorso rapido (orientamento in 1-2 sessioni):

\begin{enumerate}
\def\labelenumi{\arabic{enumi}.}
\tightlist
\item
  Capitolo 1 (problema e tesi);
\item
  Capitolo 3 (assiomi O1-O7);
\item
  Capitolo 8 (spazio di validita e triade metrica);
\item
  Capitoli 15-17 (case study + agenda scientifica).
\end{enumerate}

Percorso completo (validazione piena):

\begin{enumerate}
\def\labelenumi{\arabic{enumi}.}
\tightlist
\item
  Capitoli 1-3 (fondazione ontologica e assiomatica);
\item
  Capitoli 4-7 (estensione sistematica del corpus classico);
\item
  Capitoli 8-14 (metrica, teoremi, protocollo, benchmark, governance);
\item
  Capitoli 15-17 (casi studio e chiusura strategica).
\end{enumerate}

Regola di uso:

\begin{itemize}
\tightlist
\item
  il percorso rapido serve per capire ``cosa propone OCT'';
\item
  il percorso completo serve per valutare ``se e quanto OCT regge''.
\end{itemize}

\subsection{1.5 Stato formale degli operatori in v1.x (anti-circolarita esplicita)}

Per evitare ambiguita metodologiche, fissiamo qui lo stato formale dei tre operatori.

\begin{enumerate}
\def\labelenumi{\arabic{enumi}.}
\tightlist
\item
  Livello assiomatico (gia attivo):

  \begin{itemize}
  \tightlist
  \item
    \texttt{Coh\_Omega}, \texttt{Phi\_Omega}, \texttt{Delta\_Omega} sono definiti semanticamente e inseriti in regole decisionali esplicite.
  \end{itemize}
\item
  Livello costruttivo benchmark-specifico (attivo in forma iniziale):

  \begin{itemize}
  \tightlist
  \item
    il Capitolo 8 introduce istanziazioni calcolabili su benchmark con pipeline dichiarata.
  \end{itemize}
\item
  Livello universale cross-dominio (in consolidamento):

  \begin{itemize}
  \tightlist
  \item
    non e ancora chiuso in forma unica; resta oggetto di programma teorematico e validazione.
  \end{itemize}
\end{enumerate}

Conseguenza epistemica:

\begin{itemize}
\tightlist
\item
  in \texttt{v1.x} OCT non dichiara ``metriche universali definitive'';
\item
  dichiara invece un criterio robusto: definizione semantica + istanziazione costruttiva per dominio + stato di evidenza tracciato.
\end{itemize}

Questa scelta impedisce circolarita: la teoria non si auto-convalida per lessico, ma richiede mapping operativo e benchmark espliciti.

\begin{center}\rule{0.5\linewidth}{0.5pt}\end{center}

\section{2. Tesi del capitolo}

La tesi di questo capitolo è articolata in quattro enunciati:

\begin{enumerate}
\def\labelenumi{\arabic{enumi}.}
\tightlist
\item
  \textbf{Tesi di conservazione}: la category theory classica non va sostituita; va mantenuta come infrastruttura sintattica.
\item
  \textbf{Tesi di insufficienza}: la sola correttezza strutturale non basta per decidere la realtà operativa di una composizione.
\item
  \textbf{Tesi di estensione}: serve uno strato ordinativo esplicito definito da tre operatori minimi (\texttt{Coh\_Omega}, \texttt{Phi\_Omega}, \texttt{Delta\_Omega}).
\item
  \textbf{Tesi di scientificità}: l'estensione OCT deve restare falsificabile, con claim graduati (\texttt{S0-S3}) e protocolli di validazione riproducibili.
\end{enumerate}

Questa tesi non è una dichiarazione metaforica. È un programma tecnico: spostare la domanda da ``questa struttura è formalmente lecita?'' a ``questa struttura mantiene coerenza e genera funzione nel contesto osservato?''.

Il passaggio è minimale sul piano sintattico ma forte sul piano epistemico:

\texttt{Classical\ validity} -\textgreater{} \texttt{Classical\ validity\ +\ Ordinative\ validity}.

\begin{center}\rule{0.5\linewidth}{0.5pt}\end{center}

\section{3. Definizioni canoniche}

Per evitare ambiguità, introduciamo un lessico minimale usato in tutto il libro.

\subsection{Definizione 1 (Validità sintattica)}

Una costruzione categoriale è sintatticamente valida se soddisfa i vincoli classici del quadro in cui è definita (tipi, composizione, identità, universalità, ecc.).

Commento:

\begin{itemize}
\tightlist
\item
  questa definizione coincide con lo standard classico;
\item
  non contiene ancora alcuna ipotesi su vitalità, emergenza o degenerazione.
\end{itemize}

\subsection{Definizione 2 (Contesto osservativo)}

\texttt{Omega} è il contesto osservativo esplicito in cui viene valutata una struttura: scala, variabili, protocollo di misura, vincoli di dominio, criteri di accettazione.

Commento:

\begin{itemize}
\tightlist
\item
  OCT rifiuta valutazioni decontestualizzate;
\item
  ogni claim ontologico è indicizzato a un \texttt{Omega}.
\end{itemize}

\subsection{Definizione 3 (Coerenza ordinativa)}

\texttt{Coh\_Omega(D)} è una funzione di coerenza del diagramma/composizione \texttt{D} nel contesto \texttt{Omega}, normalizzata in \texttt{{[}0,1{]}}.

Interpretazione:

\begin{itemize}
\tightlist
\item
  valori prossimi a \texttt{1}: composizione stabile e internamente consistente;
\item
  valori prossimi a \texttt{0}: incoerenza strutturale o perdita di integrazione.
\end{itemize}

\subsection{Definizione 4 (Funzione emergente)}

\texttt{Phi\_Omega(D)} misura la funzione emergente non riducibile alla mera somma dei componenti di \texttt{D}, valutata in \texttt{Omega}.

Condizione minima:

\begin{itemize}
\tightlist
\item
  \texttt{Phi\_Omega(D)\ !=\ 0} indica presenza di funzione emergente;
\item
  \texttt{Phi\_Omega(D)\ =\ 0} indica sterilità emergente (struttura formalmente corretta ma funzionalmente piatta).
\end{itemize}

\subsection{Definizione 5 (Degenerazione ordinativa)}

\texttt{Delta\_Omega(D)\ =\ 1\ -\ Coh\_Omega(D)}.

Interpretazione:

\begin{itemize}
\tightlist
\item
  \texttt{Delta} non è un errore sintattico;
\item
  è la misura del collasso ordinativo interno.
\end{itemize}

\subsection{Definizione 6 (Validità ordinativa minima)}

Dato un valore soglia \texttt{tau\ in\ (0,1{]}}, \texttt{D} è OCT-valido in \texttt{Omega} se e solo se:

\begin{enumerate}
\def\labelenumi{\arabic{enumi}.}
\tightlist
\item
  \texttt{Coh\_Omega(D)\ \textgreater{}=\ tau}
\item
  \texttt{Phi\_Omega(D)\ !=\ 0}
\end{enumerate}

Questa è la prima forma operativa della distinzione tra ``struttura lecita'' e ``struttura reale''.

\begin{center}\rule{0.5\linewidth}{0.5pt}\end{center}

\section{4. Sviluppo formale}

\subsection{4.1 Dove il classico eccelle}

Prima di estendere, fissiamo il punto: il classico funziona in modo eccellente su:

\begin{enumerate}
\def\labelenumi{\arabic{enumi}.}
\tightlist
\item
  trasporto strutturale tra domini (funtorialità);
\item
  controllo di composizione (associatività, identità);
\item
  nozioni universali robuste (limiti, colimiti, aggiunzioni);
\item
  gerarchie di astrazione (2-categorie, arricchimenti, topos).
\end{enumerate}

In OCT, tutto questo viene preservato.

\subsection{4.2 Dove il classico è ontologicamente neutro}

La neutralità diventa criticità in tre scenari ricorrenti:

\begin{enumerate}
\def\labelenumi{\arabic{enumi}.}
\tightlist
\item
  \textbf{Equivalenza sterile}:

  \begin{itemize}
  \tightlist
  \item
    due sistemi equivalenti per struttura locale;
  \item
    esito operativo divergente su stabilità o capacità generativa.
  \end{itemize}
\item
  \textbf{Composizione degenerativa}:

  \begin{itemize}
  \tightlist
  \item
    composizione legalmente tipizzata;
  \item
    risultato interno disallineato, con degradazione progressiva.
  \end{itemize}
\item
  \textbf{Universalità vuota}:

  \begin{itemize}
  \tightlist
  \item
    limite/colimite formalmente esistente;
  \item
    nessun aumento di funzione emergente.
  \end{itemize}
\end{enumerate}

Questi tre casi non falsificano il classico. Mostrano che il classico non intendeva rispondere a quella domanda.

\subsection{4.3 Dalla ontologia estensionale alla ontologia relazionale-funzionale}

L'ontologia estensionale privilegia l'appartenenza e la sostituibilità. In domini dinamici, questa grammatica è necessaria ma non sufficiente, perché:

\begin{enumerate}
\def\labelenumi{\arabic{enumi}.}
\tightlist
\item
  la singolarità funzionale non è interamente catturata dalla classe di equivalenza;
\item
  la relazione non è solo connettivo, ma meccanismo generativo;
\item
  la funzione emergente non si deduce automaticamente dalla forma sintattica.
\end{enumerate}

OCT propone quindi un passaggio controllato:

\begin{itemize}
\tightlist
\item
  dall'oggetto come entità estensionale;
\item
  alla singolarità come nodo relazionale-funzionale.
\end{itemize}

\subsection{4.4 Principio di non-identificazione automatica}

In quadro classico:

\begin{itemize}
\tightlist
\item
  isomorfismo implica intercambiabilità strutturale rispetto alle proprietà invarianti.
\end{itemize}

In quadro OCT:

\begin{itemize}
\tightlist
\item
  isomorfismo resta valido;
\item
  ma non implica automaticamente equivalenza ordinativa piena.
\end{itemize}

Questa è una restrizione epistemica, non una negazione algebrica.

\subsection{4.5 Schema minimo di valutazione}

Dato un diagramma \texttt{D}, OCT introduce un gate in due stadi:

\begin{enumerate}
\def\labelenumi{\arabic{enumi}.}
\tightlist
\item
  gate classico: verifica sintattica;
\item
  gate ordinativo: verifica \texttt{Coh/Phi}.
\end{enumerate}

Formalmente:

\texttt{Admissible\_OCT(D,\ Omega)\ :=\ Syntactic(D)\ AND\ (Coh\_Omega(D)\ \textgreater{}=\ tau)\ AND\ (Phi\_Omega(D)\ !=\ 0)}

Il valore di \texttt{tau} è dominio-dipendente e va dichiarato nel protocollo.

\subsection{4.6 Conseguenza metodologica}

La domanda centrale cambia:

\begin{itemize}
\tightlist
\item
  classico: ``la struttura esiste?''
\item
  OCT: ``la struttura esiste e resta reale/feconda nel contesto \texttt{Omega}?''
\end{itemize}

Questo cambio è il punto di ingresso della teoria nella grammatica della scienza umana: non basta l'esistenza formale, serve una metrica di persistenza funzionale.

\subsection{4.7 Esempio astratto controllato (per chiarire la separazione)}

Consideriamo due pipeline astratte \texttt{P\_A} e \texttt{P\_B} rappresentate dallo stesso schema tipizzato:

\texttt{X\ -\textgreater{}\ Y\ -\textgreater{}\ Z}.

Nel quadro classico, se le frecce rispettano dominio/codominio e composizione, entrambe le pipeline sono formalmente ammissibili. Supponiamo ora che:

\begin{itemize}
\tightlist
\item
  in \texttt{P\_A}, la seconda freccia preservi la stabilità relazionale introdotta dalla prima;
\item
  in \texttt{P\_B}, la seconda freccia amplifichi rumore interno e riduca l'integrazione del sistema.
\end{itemize}

Dal solo punto di vista sintattico non c'è differenza decisiva: stessa forma, stesso schema compositivo, stessa legalità strutturale.

Nel quadro OCT, invece, il risultato è distinguibile:

\begin{enumerate}
\def\labelenumi{\arabic{enumi}.}
\tightlist
\item
  \texttt{Coh\_Omega(P\_A)} resta sopra soglia;
\item
  \texttt{Coh\_Omega(P\_B)} scende sotto soglia dopo iterazione;
\item
  \texttt{Phi\_Omega(P\_A)} resta non nulla;
\item
  \texttt{Phi\_Omega(P\_B)} tende a zero.
\end{enumerate}

Quindi:

\begin{itemize}
\tightlist
\item
  \texttt{P\_A} è formalmente valida e ordinativamente valida;
\item
  \texttt{P\_B} è formalmente valida ma ordinativamente invalida.
\end{itemize}

Il valore dell'esempio non è empirico, ma logico-metodologico: mostra che la differenza introdotta da OCT non è un dettaglio terminologico, ma un criterio decisionale aggiuntivo che il classico non intende fornire.

\subsection{4.8 Tre famiglie di errore inferenziale che OCT vuole evitare}

Per evitare fraintendimenti nella lettura dell'opera, fissiamo tre errori tipici che il layer ordinativo corregge.

\textbf{Errore A - Identità per sola forma}

\begin{itemize}
\tightlist
\item
  Forma: ``se due strutture sono isomorfe, allora sono anche equivalenti in esito''.
\item
  Problema: confonde equivalenza strutturale con equivalenza dinamico-funzionale.
\item
  Correzione OCT: l'isomorfismo resta, ma non chiude da solo la questione dell'esito.
\end{itemize}

\textbf{Errore B - Ottimismo compositivo automatico}

\begin{itemize}
\tightlist
\item
  Forma: ``se una composizione è ben tipizzata, allora produce progresso funzionale''.
\item
  Problema: ignora la possibilità di composizione degenerativa.
\item
  Correzione OCT: ogni composizione è sottoposta a controllo \texttt{Coh/Phi}.
\end{itemize}

\textbf{Errore C - Universalità come garanzia di fecondità}

\begin{itemize}
\tightlist
\item
  Forma: ``se esiste un limite/colimite, allora esiste anche un guadagno operativo''.
\item
  Problema: scambia proprietà universale con valore emergente.
\item
  Correzione OCT: universalità è necessaria ma non sufficiente.
\end{itemize}

Queste tre correzioni non alterano la grammatica classica. Introducono solo una disciplina inferenziale aggiuntiva quando il target non è un universo puramente formale ma un sistema con comportamento osservabile.

\subsection{4.9 Criterio di compatibilità con la teoria classica}

Per non generare una falsa antitesi, definiamo esplicitamente il criterio di compatibilità.

Compatibilità forte significa:

\begin{enumerate}
\def\labelenumi{\arabic{enumi}.}
\tightlist
\item
  ogni risultato classico resta vero nel suo dominio;
\item
  OCT non revoca teoremi classici;
\item
  OCT aggiunge condizioni di accettazione per usi ontologici/operativi.
\end{enumerate}

In pratica:

\begin{itemize}
\tightlist
\item
  se il problema è interno alla matematica pura, il layer ordinativo può essere irrilevante;
\item
  se il problema riguarda sistemi reali, il layer ordinativo diventa informativamente decisivo.
\end{itemize}

Questa distinzione evita due estremi:

\begin{enumerate}
\def\labelenumi{\arabic{enumi}.}
\tightlist
\item
  imperialismo formale (tutto è ridotto alla forma);
\item
  anti-formalismo debole (la forma non conta).
\end{enumerate}

OCT cerca una terza posizione: formalismo forte + validazione ontologica esplicita.

\subsection{4.10 Primo principio di prudenza epistemica}

Il libro adotta fin dal Capitolo 1 una regola di prudenza:

\begin{itemize}
\tightlist
\item
  nessuna inferenza sul reale senza passaggio per protocollo;
\item
  nessun protocollo senza dichiarazione di contesto;
\item
  nessun contesto senza inventario dei limiti.
\end{itemize}

Questa regola può sembrare editoriale, ma è parte della struttura teorica. Una teoria che non distingue i propri livelli di evidenza produce inevitabilmente overclaiming. L'overclaiming non è un problema retorico: è un problema strutturale, perché rende impossibile distinguere avanzamento reale da espansione narrativa.

Da qui la scelta di adottare già in questo capitolo il \texttt{Claim\ Ledger} come dispositivo di controllo interno.

\begin{center}\rule{0.5\linewidth}{0.5pt}\end{center}

\section{5. Proposizioni e teoremi locali}

\subsection{Proposizione 1.1 (Insufficienza della sola validità sintattica in domini dinamici)}

Enunciato:
in un dominio dinamico con variabili di stabilità ed emergenza osservabili, la sola validità sintattica non determina univocamente la validità operativa.

Intuizione:
due composizioni possono rispettare la stessa grammatica ma divergere in robustezza interna.

Stato:

\begin{itemize}
\tightlist
\item
  claim teorico forte, da validare per classi di dominio (\texttt{S1/S2} a seconda del dataset).
\end{itemize}

\subsection{Proposizione 1.2 (Necessità del contesto osservativo)}

Enunciato:
ogni valutazione di validità ordinativa è indicizzata da un contesto \texttt{Omega}; non esiste valore assoluto decontestualizzato di \texttt{Coh} e \texttt{Phi}.

Intuizione:
la stessa architettura può risultare coerente a una scala e degenerativa a un'altra.

Stato:

\begin{itemize}
\tightlist
\item
  principio metodologico (\texttt{S1}).
\end{itemize}

\subsection{Teorema locale 1.3 (Separazione dei livelli di validità)}

Ipotesi:

\begin{enumerate}
\def\labelenumi{\arabic{enumi}.}
\tightlist
\item
  \texttt{D} è sintatticamente valido in una categoria \texttt{C}.
\item
  Esiste un contesto \texttt{Omega} in cui \texttt{Coh\_Omega(D)\ \textless{}\ tau} oppure \texttt{Phi\_Omega(D)\ =\ 0}.
\end{enumerate}

Tesi:
\texttt{D} è formalmente ammissibile ma non OCT-valido in \texttt{Omega}.

Proof sketch:
dal punto 1 segue l'ammissibilità classica.
dal punto 2 segue la violazione della definizione di validità ordinativa minima.
la congiunzione richiede entrambi i gate; quindi il gate ordinativo fallisce.

Conseguenze:

\begin{enumerate}
\def\labelenumi{\arabic{enumi}.}
\tightlist
\item
  la validità classica è necessaria ma non sufficiente;
\item
  OCT non contraddice il classico, lo rifinisce con una condizione additiva.
\end{enumerate}

Stato:

\begin{itemize}
\tightlist
\item
  teorema locale definizionale (\texttt{in\_proof}, orientato a \texttt{S1} su casi testati).
\end{itemize}

\subsection{Corollario 1.4 (Non-ridondanza della OCT rispetto alla teoria classica)}

Se esiste almeno una famiglia non vuota di diagrammi sintatticamente validi ma ordinativamente invalidi, allora l'estensione OCT non è ridondante rispetto al quadro classico.

Commento:
questo corollario definisce la condizione minima di senso dell'intero progetto.

\begin{center}\rule{0.5\linewidth}{0.5pt}\end{center}

\section{6. Implicazioni operative}

Le implicazioni immediate del capitolo sono sei.

\begin{enumerate}
\def\labelenumi{\arabic{enumi}.}
\item
  \textbf{Design teorico}

  \begin{itemize}
  \tightlist
  \item
    ogni capitolo successivo deve esplicitare cosa conserva del classico e cosa filtra/estende.
  \end{itemize}
\item
  \textbf{Design sperimentale}

  \begin{itemize}
  \tightlist
  \item
    ogni caso studio deve dichiarare \texttt{Omega}, metrica \texttt{Coh}, criterio di \texttt{Phi}.
  \end{itemize}
\item
  \textbf{Design editoriale}

  \begin{itemize}
  \tightlist
  \item
    nessun enunciato resta in forma ``assoluta'': va sempre indicato livello di evidenza.
  \end{itemize}
\item
  \textbf{Design comparativo}

  \begin{itemize}
  \tightlist
  \item
    confronto classico/OCT effettuato su coppie di casi:

    \begin{enumerate}
    \def\labelenumii{\arabic{enumii}.}
    \tightlist
    \item
      stesso statuto sintattico;
    \item
      differenza su coerenza/emergenza.
    \end{enumerate}
  \end{itemize}
\item
  \textbf{Design teorematico}

  \begin{itemize}
  \tightlist
  \item
    teoremi futuri devono specificare se sono:

    \begin{enumerate}
    \def\labelenumii{\arabic{enumii}.}
    \tightlist
    \item
      conservativi (traslano risultati classici),
    \item
      differenziali (impossibili senza layer ordinativo),
    \item
      applicativi (testabili su benchmark).
    \end{enumerate}
  \end{itemize}
\item
  \textbf{Design di governance scientifica}

  \begin{itemize}
  \tightlist
  \item
    i claim non validati restano \texttt{revise};
  \item
    la teoria cresce per cicli di validazione, non per accumulo retorico.
  \end{itemize}
\end{enumerate}

\subsection{6.1 Protocollo minimo applicabile già dal Capitolo 2}

Per rendere il lavoro cumulativo, introduciamo un protocollo minimo in séi passi che verrà riusato nei capitoli successivi.

\begin{enumerate}
\def\labelenumi{\arabic{enumi}.}
\tightlist
\item
  Definire l'unità strutturale analizzata (\texttt{D}).
\item
  Dichiarare il contesto osservativo \texttt{Omega}.
\item
  Esplicitare la soglia \texttt{tau} con motivazione.
\item
  Stimare \texttt{Coh\_Omega(D)} con metrica dichiarata.
\item
  Stimare \texttt{Phi\_Omega(D)} con criterio di non-sterilità.
\item
  Classificare esito:

  \begin{itemize}
  \tightlist
  \item
    \texttt{accepted} se passa entrambi i gate;
  \item
    \texttt{rework} se passa il gate sintattico ma fallisce quello ordinativo;
  \item
    \texttt{rejected} se fallisce anche i vincoli formali.
  \end{itemize}
\end{enumerate}

Il protocollo non sostituisce i metodi dominio-specifici. Serve a garantire che ogni estensione OCT sia confrontabile e auditabile.

\subsection{6.2 Valore operativo per la stesura dell'intero libro}

Il Capitolo 1 impone una disciplina che influenza direttamente la scrittura del testo completo:

\begin{enumerate}
\def\labelenumi{\arabic{enumi}.}
\tightlist
\item
  ogni capitolo dovrà distinguere ``estensione definizionale'' da ``validazione effettiva'';
\item
  ogni caso studio dovrà dichiarare cosa è \texttt{S0}, cosa è \texttt{S1}, cosa resta \texttt{S2/S3};
\item
  ogni proposta forte dovrà includere almeno un failure mode plausibile.
\end{enumerate}

In assenza di questa disciplina, il libro rischierebbe di oscillare tra manifesto teorico e report tecnico senza stabilire una grammatica unica. Con questa disciplina, invece, il testo può mantenere una voce scientifica coerente lungo tutte le 50.000 parole.

\begin{center}\rule{0.5\linewidth}{0.5pt}\end{center}

\section{7. Limiti e non-claim}

Questo capitolo non dimostra:

\begin{enumerate}
\def\labelenumi{\arabic{enumi}.}
\tightlist
\item
  che \texttt{Coh\_Omega} abbia già una metrica unica universale valida per tutti i domini;
\item
  che \texttt{Phi\_Omega} sia già calibrata in modo definitivo nei casi sociali/cognitivi;
\item
  che ogni divergenza empirica tra sistemi isomorfi richieda necessariamente un framework ordinativo;
\item
  che la sola introduzione di \texttt{Coh/Phi/Delta} risolva automaticamente i problemi di causalità complessa;
\item
  che i casi studio già disponibili abbiano status \texttt{validated} definitivo su tutta la pipeline.
\end{enumerate}

Ulteriore limite:

\begin{itemize}
\tightlist
\item
  qui è stata definita una soglia minima di validità ordinativa (\texttt{tau}), ma non una teoria completa di scelta ottima di \texttt{tau} per dominio. Questa scelta resta metodologica e va motivata in capitoli successivi.
\end{itemize}

Non-claim esplicito:

\begin{itemize}
\tightlist
\item
  OCT non è presentata come alternativa antagonista alla category theory classica;
\item
  OCT è presentata come estensione conservativa con gate ontologico aggiuntivo.
\end{itemize}

\begin{center}\rule{0.5\linewidth}{0.5pt}\end{center}

\section{8. Claim Ledger (S0-S3)}

{\def\LTcaptype{none} % do not increment counter
\begin{longtable}[]{@{}lllll@{}}
\toprule\noalign{}
ID & Claim & Tipo & Evidenza & Stato \\
\midrule\noalign{}
\endhead
\bottomrule\noalign{}
\endlastfoot
C1.1 & La validità sintattica da sola è insufficiente in domini dinamici & teorico-metodologico & S1 & revise \\
C1.2 & Serve indicizzare la validità a un contesto osservativo \texttt{Omega} & metodologico & S1 & candidate \\
C1.3 & La distinzione classico/OCT può essere formalizzata come doppio gate & formale & S1 & in\_proof \\
C1.4 & L'estensione OCT è non ridondante se esistono casi classically-valid ma OCT-invalid & metateorico & S2 & revise \\
C1.5 & \texttt{Coh}, \texttt{Phi}, \texttt{Delta} costituiscono la triade minima di validazione ordinativa & architetturale & S2 & candidate \\
C1.6 & In \texttt{v1.x} la triade e usabile senza circolarita solo se accompagnata da istanziazioni costruttive benchmark-specifiche & metodologico & S1 & candidate \\
\end{longtable}
}

Note:

\begin{itemize}
\tightlist
\item
  in questo capitolo \texttt{S0} non compare, perché non sono introdotti nuovi dataset empirici;
\item
  la promozione a \texttt{validated} richiede benchmark espliciti nei capitoli metodologici.
\end{itemize}

\begin{center}\rule{0.5\linewidth}{0.5pt}\end{center}

\section{9. Bridge al Capitolo 2}

Capitolo 1 ha introdotto il problema e il criterio minimo di validità ordinativa.
Capitolo 2 affronta la transizione ontologica completa: dalla ``cosa'' alla ``singolarità''.

Obiettivo del Capitolo 2:

\begin{enumerate}
\def\labelenumi{\arabic{enumi}.}
\tightlist
\item
  mostrare perché l'identità ordinativa non collassa nella sola equivalenza estensionale;
\item
  formalizzare la nozione di singolarità funzionale;
\item
  derivare il primo set di invarianti relazionali che preparano la definizione piena di categoria ordinativa nel Capitolo 3.
\end{enumerate}

In termini operativi:
se il Capitolo 1 ha definito \textbf{quando} una composizione è valida,
il Capitolo 2 definira \textbf{che cosa} viene composto in OCT.

\begin{center}\rule{0.5\linewidth}{0.5pt}\end{center}

\section{Riferimenti}

\begin{enumerate}
\def\labelenumi{\arabic{enumi}.}
\tightlist
\item
  Ghioni, F. (2026). \emph{TE\_CORE v5.1: Technology of Expressions - Core Ontology and Axioms}.
\item
  Ghioni, F. (2026). \emph{Ordinative Set Theory (OST): Concise Operational Guide for AI, v2.1}.
\item
  Ghioni, F. (2026). \emph{OCT Typed Formal Specification v0.1}.
\item
  Mac Lane, S. (1998). \emph{Categories for the Working Mathematician} (2nd ed.).
\item
  Riehl, E. (2016). \emph{Category Theory in Context}.
\item
  Awodey, S. (2010). \emph{Category Theory} (2nd ed.).
\end{enumerate}
